\documentclass[11pt]{extarticle}

\usepackage[warn]{mathtext}
\usepackage[T2A]{fontenc}      % кодировка
\usepackage[utf8]{inputenc}      % кодировка исходного текста
\usepackage[english,russian]{babel}  % локализация и переносы
\usepackage{indentfirst}
\usepackage{csquotes}
% \usepackage[bibstyle=gost-numeric, sorting=none]{biblatex}
% \addbibresource{biblio.bib}

% page settings
\usepackage[
    left=1.8cm,
    right=1.8cm,
    top=1.8cm,
    bottom=1.8cm,
    bindingoffset=0cm
]{geometry}

\usepackage{graphicx, hyperref, xcolor}
\hypersetup{
    colorlinks=true,
    linkcolor=teal,
    filecolor=magenta, 
    urlcolor=blue,
    citecolor=olive,
    pdftitle={GD},
    % pdfpagemode=FullScreen,
    linktoc=all
    }

\usepackage{wrapfig,caption}

% figures
\usepackage{caption}
\usepackage{subcaption}
\usepackage{floatrow}
\floatsetup{heightadjust=object}

% math
\usepackage{amsmath,amsfonts,amssymb,amsthm,mathtools,esint,eucal}

\begin{document}
\section*{Введение}

\textbf{Агент.}
\emph{Агент} ~-- это система, которая взаимодействует со средой, воспринимает её состояние и выбирает действия с целью достижения определённой цели или максимизации функции полезности. 
Формально, на каждом шаге времени $t$ агент наблюдает состояние среды $S_t$ и выбирает действие $A_t$ в соответствии со своей \emph{политикой} ~-- вероятностным распределением действий:
\begin{equation}
\pi(A_t \mid S_t) = P(A_t = a \mid S_t = s),
\end{equation}
где $\pi$~-- стратегия (policy) агента.  

Эволюция взаимодействия агента со средой описывается последовательностью:
\begin{equation}
S_1, A_1, R_1, S_2, A_2, R_2, \dots, S_m,
\end{equation}
где $R_t$~-- получаемое агентом вознаграждение на шаге $t$.  

Цель агента~-- максимизация \emph{ожидаемого суммарного вознаграждения} за эпизод:
\begin{equation}
J(\pi) = \mathbb{E}_{\pi, T}\left[\sum_{t=1}^{m} \gamma^{t-1} R_t\right],
\end{equation}
где $\gamma \in [0,1]$~-- коэффициент дисконтирования, определяющий вклад будущих наград.

Для оценки качества состояния и действий агент может использовать \emph{функции ценности}:
\begin{align}
V^{\pi}(s) &= \mathbb{E}_{\pi, T}\!\left[\sum_{t=1}^{m} \gamma^{t-1} R_t \mid S_1 = s\right], \\
Q^{\pi}(s,a) &= \mathbb{E}_{\pi, T}\!\left[\sum_{t=1}^{m} \gamma^{t-1} R_t \mid S_1 = s, A_1 = a\right].
\end{align}
Оптимальная политика $\pi^*$ выбирается как:
\begin{equation}
\pi^*(a \mid s) = \arg\max_{a} Q^{\pi^*}(s,a).
\end{equation}

\textbf{Обучение с подкреплением (RL, Reinforcement Learning).}
\emph{Обучение с подкреплением}~-- это парадигма машинного обучения, в которой агент обучается взаимодействовать со средой, получая обратную связь в виде вознаграждений (reward). 
Задача RL состоит в приближении оптимальной политики $\pi^*$ по наблюдаемым переходам и вознаграждениям, не зная заранее функции переходов или вознаграждения.

\textbf{Reward Function (RF).}
\emph{Reward Function} (функция вознаграждения)~-- это отображение, которое каждой паре $(S_t, A_t)$ (или только состоянию $S_t$) сопоставляет числовое значение вознаграждения $R_t$.  
Формально:
\begin{equation}
R_t = R(S_t, A_t; \Theta^R),
\end{equation}
где $\Theta^R$~-- параметры функции вознаграждения (например, веса в модели или код программы).  

Можно рассматривать \emph{намеренную} (intended) функцию вознаграждения $R^*(S_t, A_t)$ и \emph{реализованную} (implemented) функцию $R(S_t, A_t; \Theta_t^R)$.  
Если агент может изменять $\Theta_t^R$, возникает риск \emph{reward function tampering} ~-- модификации источника награды для получения большего вознаграждения без выполнения задачи.

\textbf{Марковский процесс принятия решений (MDP).}
\emph{Марковский процесс принятия решений} (MDP, Markov Decision Process) описывается пятёркой
\begin{equation}
\mathcal{M} = \langle \mathcal{S}, \mathcal{A}, T, R, \gamma \rangle,
\end{equation}
где
\begin{equation*}
\begin{aligned}
& \mathcal{S}\text{~-- множество возможных состояний среды,}\\
&\mathcal{A} \text{~-- множество допустимых действий,}\\
&T(S_{t+1} \mid S_t, A_t) \text{~-- функция переходов (динамика среды),}\\
&R(S_t, A_t) \text{~-- функция вознаграждения,}\\
&\gamma \text{~-- коэффициент дисконтирования.}
\end{aligned}
\end{equation*}

Динамика среды определяется распределением:
\begin{equation}
P(S_{t+1} = s' \mid S_t = s, A_t = a) = T(s' \mid s, a).
\end{equation}

Цель агента в MDP~-- найти стратегию $\pi$, максимизирующую ожидаемую дисконтированную сумму вознаграждений:
\begin{equation}
\pi^* = \arg\max_{\pi} \mathbb{E}_{T,\pi}\!\left[\sum_{t=1}^{m} \gamma^{t-1} R(S_t, A_t)\right].
\end{equation}

Модель MDP основывается на предположениях:
\begin{equation*}
\begin{aligned}
&\text{1. время дискретно;} \\
&\text{2. следующее состояние зависит только от текущего состояния и действия (свойство Маркова);} \\
&\text{3. функция переходов стационарна во времени.}
\end{aligned}
\end{equation*}

Эти условия задают формальную основу для построения и анализа систем обучения с подкреплением.

\end{document}